\documentclass[11pt]{article}
\usepackage[margin=1in]{geometry}
\usepackage{amsmath,amssymb,booktabs,hyperref}
\title{Replication Report: Field-Free Deterministic SOT Switching via\\
Cluster-Octupole Out-of-Plane Spin Polarization $\sigma_z$ (You et al., 2021)}
\author{Automated replication (macrospin LLG)}
\date{\today}
\begin{document}
\maketitle

\section{Target claim}
You \emph{et al.} (2021, Mn$_3$SnN antiperovskite noncollinear AFM) report that an
out-of-plane spin polarization $\boldsymbol{\sigma}_z\propto\mathbf{H}_{so}\times\mathbf{T}$,
generated by the cluster magnetic octupole $\mathbf{T}$ when the current $\mathbf{J}\parallel\mathbf{T}$,
drives \textbf{field-free deterministic} spin--orbit-torque (SOT) switching of an adjacent
perpendicular ferromagnet (Co/Pd)$_3$. The switching polarity is set by the current sign,
and $\sigma_z$ (hence the deterministic field-free switching) \emph{vanishes} when
$\mathbf{J}\perp\mathbf{T}$. Measured torque ratios: $\theta_{AD,z}=0.003$, $\theta_{FL,z}=0.053$;
$J_{crit}\approx 9\times10^{6}\,\mathrm{A\,cm^{-2}}$.

\section{Nature of the paper and scope of replication}
The paper is predominantly \emph{experimental} (epitaxial growth, XRD, SQUID, AHE,
ST-FMR, MOKE, transport). No simulation workflow or code is provided. However, its
central mechanistic assertion is a \emph{symmetry-defined theory result}: the antidamping
$\sigma_z$ torque $\tau_C=\mathbf{m}\times(\mathbf{m}\times\boldsymbol{\sigma}_z)$ lifts the
$\pm\hat z$ degeneracy of a PMA magnet with no external field. We reproduce this
from scratch with a single-macrospin Landau--Lifshitz--Gilbert (LLG) integration.
The materials-growth, ST-FMR fitting, and MOKE imaging portions are \emph{not}
reproducible without the physical samples and instruments; hence the overall verdict is
\textbf{PARTIAL}.

\section{Model}
Unit magnetization $\mathbf{m}$, uniaxial PMA along $\hat z$
($\mathbf{H}=h_k m_z\hat z$, $h_k=1$), Gilbert damping $\alpha=0.1$, reduced units,
\textbf{zero external field everywhere}. RK4 integration of
\[
\frac{d\mathbf{m}}{dt}=-\mathbf{m}\times\mathbf{H}
-\alpha\,\mathbf{m}\times(\mathbf{m}\times\mathbf{H})
+\tau_{DL}\,\mathbf{m}\times(\mathbf{m}\times\mathbf{p})
+\tau_{FL}\,\mathbf{m}\times\mathbf{p},
\]
with SOT amplitudes scaled from the paper's ratios, $\tau_{DL}=C\theta_{AD,z}$,
$\tau_{FL}=C\theta_{FL,z}$, current pulse ON for 60\% of the run then relaxed.
\begin{itemize}
\item \textbf{Case A} ($\mathbf{J}\parallel\mathbf{T}$, $\sigma_z$ present): polarization
$\mathbf{p}\approx\hat z$ (dominant out-of-plane, small $\sigma_y$).
\item \textbf{Case B} ($\mathbf{J}\perp\mathbf{T}$, $\sigma_z$ absent): $\mathbf{p}=\hat y$
(conventional in-plane $\sigma_y$ only).
\end{itemize}
Determinism test: initialise both $m_z=+1$ and $m_z=-1$; require the final state to depend
\emph{only} on current sign, not on initial state.

\section{Results}
\begin{table}[h]\centering
\begin{tabular}{llcc}
\toprule
Case & init & final ($+I$) & final ($-I$)\\
\midrule
A ($\sigma_z$ present) & up   & down & up\\
A ($\sigma_z$ present) & down & down & up\\
B ($\sigma_z$ absent)  & up   & up   & up\\
B ($\sigma_z$ absent)  & down & down & down\\
\bottomrule
\end{tabular}
\caption{Case A: final state set by current sign, independent of initial state
$\Rightarrow$ deterministic field-free switching. Case B: final state tracks initial state,
no switching. Threshold current $C_{crit}\approx 12$ (reduced units).}
\end{table}

Case A reproduces \textbf{deterministic field-free switching with current-sign-controlled
polarity}. Case B shows \textbf{no} deterministic field-free switching. This matches the
paper's key control experiment (Fig.~4c vs.~4f: switching along [001] $J\parallel T$,
no switching along [110] $J\perp T$).

\section{Verdict: PARTIAL}
\textbf{Coverage 6/10} --- the core mechanistic theory claim is fully reproduced; the
experimental growth/ST-FMR/MOKE bulk of the paper is not reproducible here.
\textbf{Agreement 8/10} --- qualitative and structural agreement is complete
(deterministic switching present iff $\sigma_z$ present; polarity set by current sign;
finite threshold current). Absolute $J_{crit}$ in physical units and the small
$\theta_{AD,z}$ are not quantitatively benchmarked, only used as the FL/AD ratio.

\section{Physics summary}
The antidamping component of the out-of-plane spin polarization,
$\tau_C\propto\mathbf{m}\times(\mathbf{m}\times\boldsymbol{\sigma}_z)$, is the only torque
term that can deterministically break the up/down symmetry of a perpendicular magnet
without an external field. Our macrospin LLG confirms this: with $\sigma_z$ present the
final magnetization is fixed by the current sign regardless of the starting state, and with
$\sigma_z$ removed (only $\sigma_y$) the magnet is not switched --- exactly the
$J\parallel T$ vs.\ $J\perp T$ dichotomy the paper attributes to the cluster magnetic octupole.
\end{document}
